\section{Extensions and open items}\label{sec:sp:extensions}

This appendix records two different kinds of unfinished business. Section~\ref{sec:sp:ext:open} lists
gaps in the derivation as it stands --- things that ought to be established for the present model and
have not been. Section~\ref{sec:sp:ext:later} lists extensions to the model itself: modelling choices
whose alternatives are described but not derived anywhere in this note, and which are candidates for a
separate project rather than for this paper. The two lists are kept apart deliberately, because they
carry different burdens: the first is work owed, the second is work optional.


\subsection{Open items in the present derivation}\label{sec:sp:ext:open}

\begin{enumerate}[label=\textbf{O\arabic*.},ref=O\arabic*,leftmargin=2.6em,itemsep=0.35em]

\item\label{op:sp:soc} \textbf{Sufficiency and uniqueness.} The conditions of
Sections~\ref{sec:sp:A} and~\ref{sec:sp:B} are necessary only. Concavity of the maximised Hamiltonian
in the states is not established: $F$, $u$ and $-v$ are concave, but the ledger is bilinear in $\varpi$
and $W$, the recycling technology enters through a constant-returns function of two controls, and under
Version A the intensity definition is bilinear in $\Omega$ and the flows in $\mathcal D$. Uniqueness is
not established either, and it is invoked when the market economy is said to implement \emph{the}
planner's allocation. See Remarks~\ref{rem:sp:A:soc} and~\ref{rem:sp:B:soc}.

\item\label{op:sp:corner} \textbf{Behaviour at the corners.} The treatment margin was shown never to
sit at $\varpi=0$ (Sections~\ref{sec:sp:A:foc} and~\ref{sec:sp:B:foc}), but the $\varpi=1$ corner, the
$D=0$ corner, and the corner $\eta>0$ where the Version-B mass constraint~\eqref{eq:sp:B:mass} binds
have not been characterised. The last of these matters for the numerical solution, which must handle a
complementarity condition rather than a smooth system --- and it is a reason to prefer Version A for a
first numerical pass, since the constraint is vacuous there.

\item\label{op:sp:exist} \textbf{Existence of an optimum, and of the asymptotic path.}
Propositions~\ref{prop:sp:A:nosteadystate} and~\ref{prop:sp:B:nosteadystate} rule out an interior rest
point, and Sections~\ref{sec:sp:A:longrun} and~\ref{sec:sp:B:longrun} argue informally that the economy
approaches a materially closed limit. That the limit is attained, that the transversality
conditions~\eqref{eq:sp:B:tvc} hold along it, and that lifetime utility is finite, are all unverified.
Switch~\ref{sw:sp:CDX} sidesteps the issue by restoring a steady state; it is not a proof that the
general case is well posed.

\item\label{op:sp:period} \textbf{Sensitivity to the recirculation assumption.} Recycled material
re-enters production contemporaneously, $R^R=a\varpi W$ at the same instant. In continuous time this
is innocuous; in a discrete-time implementation it makes the circularity multiplier
$\left(1-a\varpi\right)^{-1}$ a function of the period length, and the multiplier diverges in exactly
the region the paper is about. The ratios $R/N$ and $W/N$ should be reported alongside every headline
result, and their sensitivity to the period length checked.

\item\label{op:sp:nonopt} \textbf{Welfare under non-optimal policy.} Nothing here characterises the
allocation, or $U_0$, under an arbitrary policy that is not the first best. Any counterfactual
exercise that compares the optimum to a status-quo or partial-policy scenario needs this, and it does
not follow from the conditions above.

\item\label{op:sp:identify} \textbf{What the calibration can identify.} Under Version B the relative
intensities $\omega^j$ do not affect the allocation
(Propositions~\ref{prop:sp:B:irrelevance} and~\ref{prop:sp:inert}), so no feature of the optimal path
carries information about them; under Version A they do. The calibration strategy therefore has to be
settled jointly with the choice of closure, and it is not settled here. See
Section~\ref{sec:sp:B:compare}.

\end{enumerate}


\subsection{Extensions for later work}\label{sec:sp:ext:later}

\begin{enumerate}[label=\textbf{E\arabic*.},ref=E\arabic*,leftmargin=2.6em,itemsep=0.35em]

\item\label{ext:sp:manna} \textbf{A stock for nature-attributable mass.} The flows
$\Omega^{N,S}N$ and $\Omega^{D,S}D$ enter the waste ledger without being drawn from any stock ---
overburden arrives from nowhere. Section~\ref{sec:sp:setup:materialaccounting} flags this and assumes
it is not the binding constraint. Closing it would require a second natural stock with its own law of
motion, and would make the mass balance of the \emph{whole} system, not just of the economy, hold
exactly.

\item\label{ext:sp:lag} \textbf{A recycling lag.} Replacing $R^R_t=a\varpi W_t$ by a distributed lag
over past waste generation adds at least one state (the pool of waste awaiting treatment) and removes
the closed-form solution of the ledger. This is the single maintained assumption most likely to matter
quantitatively; see~\ref{op:sp:period}.

\item\label{ext:sp:pools} \textbf{Two embodied-material pools.} The note carries one well-mixed
stock $M^K$, so $K^Y$ and $K^R$ share the average intensity $\Phi^K$. Distinguishing them requires a
second state and a second liability price, and would matter whenever the two sectors' investment
histories differ --- which they do exactly during the transition, when $K^R$ is growing fastest.

\item\label{ext:sp:reuse} \textbf{Direct reuse of depreciated capital material.} Depreciation
currently returns embodied material to the undifferentiated waste flow, where it must be collected and
treated like anything else. An alternative is costless, or cheaper, re-embodiment of a share of it ---
remanufacturing rather than recycling. This changes the ledger and requires $q$ to be re-derived, since
capital would then be a partially permanent rather than purely temporary sink.

\item\label{ext:sp:residual} \textbf{Residual waste from recycling as a separate flow.} The recycling
function is interpreted as net of the treatment of its own residuals
(Section~\ref{sec:sp:setup:materialaccounting}). Making the residual explicit requires a second waste
stream with its own $d^W$, and only earns its keep in a model with more than one material.

\item\label{ext:sp:multi} \textbf{Multiple materials.} A single homogeneous tonne is the strongest
simplification in the note. Distinct materials differ in recyclability $a(\cdot)$, in the emission
share $d^W$, in extraction cost curvature, and in substitutability inside $F$; and the composition of
the material bundle is itself an endogenous margin that the present model cannot speak to. This is the
most substantial of the extensions listed here and is properly a separate paper.

\item\label{ext:sp:landfill} \textbf{Landfill as a stock.} Untreated waste currently emits a fraction
$1$ and treated waste a fraction $d^W$, both instantaneously. Storing waste in a stock that leaks
slowly, and that can be mined later, would connect the model to the literature on legacy waste and
would give a second, non-natural reserve from which material can be drawn once $p^S$ is high enough.

\item\label{ext:sp:Aackstop} \textbf{A bounded resource base or a backstop.} Exploration is
unbounded, so the resource is never exhausted, only ever more expensive
(Section~\ref{sec:sp:setup:primitives}); this is what
Propositions~\ref{prop:sp:A:nosteadystate} and~\ref{prop:sp:B:nosteadystate} turn on. Imposing an upper bound on cumulative discovery, or
a backstop technology that caps $C^N_N$, would change the long-run character of the model
substantially and should be compared against switch~\ref{sw:sp:CDX}.

\item\label{ext:sp:tech} \textbf{Endogenous recycling technology.} $a(\cdot)$, $c^T(\cdot)$ and
$\Phi^I$ are exogenous functions. Making any of them respond to R\&D or to learning-by-doing would add
a state and would give the transition a second driver alongside scarcity --- and, unlike scarcity, one
that policy can act on directly.

\item\label{ext:sp:design} \textbf{Product design: endogenous $\delta$ and $\Phi^I$.} Durability and
material content per unit of capital are choices, not parameters. Endogenising them is the natural
formalisation of ``design for circularity'' and would make $q$ a genuine two-margin object. Under
Version A, $\Phi^I$ is already partly endogenous, but through the accounting rather than through a
design decision.

\item\label{ext:sp:open} \textbf{Open economy.} Imported virgin material, exported waste, and a world
price for both. This changes what the resource rent means for a single country and is the margin most
of the applied circular-economy policy debate is actually about.

\item\label{ext:sp:hetero} \textbf{Heterogeneity and distribution.} A representative dynasty rules out
any question about who bears the cost of the transition and who receives the instrument revenues,
which is where most of the political economy of these policies sits.

\item\label{ext:sp:uncertainty} \textbf{Uncertainty.} Reserves, damages and recycling technology are
all deterministic. Reserve uncertainty in particular interacts with exploration in a way the current
specification cannot represent, since exploration here reduces cost rather than resolving risk.

\end{enumerate}
